\documentclass[11pt]{article}
% ===================================================================
% Shared preamble for the "tiny lemma, read slowly" preprint series.
% Mirrors the style of FrobeniusLemmas_preprint.tex.
% ===================================================================

\usepackage{amsmath,amssymb,amsthm}
\usepackage{mathtools}
\usepackage[margin=1.2in]{geometry}
\usepackage{hyperref}
\usepackage{xcolor}
\usepackage{listings}
\usepackage{tikz}
\usepackage{tikz-cd}
\usepackage{mathpartir}
\usepackage{newunicodechar}
\usepackage{setspace}

\onehalfspacing

% ---- Unicode glyphs ----
\newunicodechar{→}{\ensuremath{\to}}
\newunicodechar{↦}{\ensuremath{\mapsto}}
\newunicodechar{↔}{\ensuremath{\leftrightarrow}}
\newunicodechar{⟹}{\ensuremath{\Longrightarrow}}
\newunicodechar{⟶}{\ensuremath{\longrightarrow}}
\newunicodechar{≠}{\ensuremath{\neq}}
\newunicodechar{≤}{\ensuremath{\leq}}
\newunicodechar{≥}{\ensuremath{\geq}}
\newunicodechar{∀}{\ensuremath{\forall}}
\newunicodechar{∃}{\ensuremath{\exists}}
\newunicodechar{∘}{\ensuremath{\circ}}
\newunicodechar{·}{\ensuremath{\cdot}}
\newunicodechar{∈}{\ensuremath{\in}}
\newunicodechar{∉}{\ensuremath{\notin}}
\newunicodechar{≃}{\ensuremath{\simeq}}
\newunicodechar{≅}{\ensuremath{\cong}}
\newunicodechar{≈}{\ensuremath{\approx}}
\newunicodechar{∧}{\ensuremath{\wedge}}
\newunicodechar{∨}{\ensuremath{\vee}}
\newunicodechar{¬}{\ensuremath{\neg}}
\newunicodechar{⊢}{\ensuremath{\vdash}}
\newunicodechar{⟨}{\ensuremath{\langle}}
\newunicodechar{⟩}{\ensuremath{\rangle}}
\newunicodechar{⇑}{\ensuremath{\mathord{\uparrow}}}
\newunicodechar{↑}{\ensuremath{\mathord{\uparrow}}}
\newunicodechar{∎}{\ensuremath{\blacksquare}}
\newunicodechar{∣}{\ensuremath{\mid}}
\newunicodechar{∤}{\ensuremath{\nmid}}
\newunicodechar{∑}{\ensuremath{\textstyle\sum}}
\newunicodechar{∏}{\ensuremath{\textstyle\prod}}
\newunicodechar{∅}{\ensuremath{\emptyset}}
\newunicodechar{∪}{\ensuremath{\cup}}
\newunicodechar{∩}{\ensuremath{\cap}}
\newunicodechar{≡}{\ensuremath{\equiv}}
\newunicodechar{ℕ}{\ensuremath{\mathbb{N}}}
\newunicodechar{ℤ}{\ensuremath{\mathbb{Z}}}
\newunicodechar{ℝ}{\ensuremath{\mathbb{R}}}
\newunicodechar{𝔽}{\ensuremath{\mathbb{F}}}
\newunicodechar{α}{\ensuremath{\alpha}}
\newunicodechar{β}{\ensuremath{\beta}}
\newunicodechar{σ}{\ensuremath{\sigma}}
\newunicodechar{τ}{\ensuremath{\tau}}
\newunicodechar{φ}{\ensuremath{\varphi}}
\newunicodechar{ψ}{\ensuremath{\psi}}
\newunicodechar{χ}{\ensuremath{\chi}}
\newunicodechar{Φ}{\ensuremath{\Phi}}
\newunicodechar{Δ}{\ensuremath{\Delta}}
\newunicodechar{₀}{\textsubscript{0}}
\newunicodechar{₁}{\textsubscript{1}}
\newunicodechar{₂}{\textsubscript{2}}
\newunicodechar{ᵏ}{\ensuremath{{}^{k}}}
\newunicodechar{ⁿ}{\ensuremath{{}^{n}}}
\newunicodechar{ʲ}{\ensuremath{{}^{j}}}
\newunicodechar{²}{\ensuremath{{}^{2}}}
\newunicodechar{³}{\ensuremath{{}^{3}}}
\newunicodechar{⁴}{\ensuremath{{}^{4}}}
\newunicodechar{⁻}{\ensuremath{{}^{-}}}
\newunicodechar{¹}{\ensuremath{{}^{1}}}

% ---- Lean listing style ----
\definecolor{leanbg}{RGB}{247,247,250}
\definecolor{leankw}{RGB}{0,90,160}
\definecolor{leancm}{RGB}{120,120,120}
\lstdefinestyle{lean}{
  backgroundcolor=\color{leanbg},
  basicstyle=\ttfamily\small,
  keywordstyle=\color{leankw}\bfseries,
  commentstyle=\color{leancm}\itshape,
  morekeywords={theorem,lemma,def,fun,Prop,Type,variable,where,
    Field,Fintype,Finite,DecidableEq,CommRing,CommSemiring,Ring,CharP,Function,
    Injective,Surjective,Bijective,by,rfl,have,rw,exact,ext,simp,intro,
    iterateFrobenius,RingHom,Commute,convert,apply,using,Nat,Coprime,gcd,
    Odd,Even,IsAPN,IsAB,walsh,Nat,obtain,induction,cases,calc,grind},
  showstringspaces=false,
  columns=fullflexible,
  extendedchars=true,
  frame=single,
  framesep=6pt,
  xleftmargin=4pt,
  aboveskip=12pt,
  belowskip=14pt,
}
\lstset{style=lean}

\newtheorem{lemma}{Lemma}
\newtheorem{theorem}{Theorem}
\newtheorem{corollary}{Corollary}

% boxed "what just happened" note
\newcommand{\note}[1]{%
  \par\smallskip
  \noindent\colorbox{leanbg}{\parbox{\dimexpr\linewidth-2\fboxsep}{\small #1}}%
  \par\smallskip}


\title{\textbf{One gcd, four floors up}\\[4pt]
\large $\gcd(2^a-1,\,2^b-1)=2^{\gcd(a,b)}-1$, read slowly}
\author{}
\date{}

\begin{document}
\maketitle

\begin{abstract}
\noindent
One fact about integers.\\[2pt]
$\bullet$\quad \texttt{gcd\_mersenne}:\;\; $\gcd(2^a-1,\,2^b-1)=2^{\gcd(a,b)}-1$.\\[6pt]
\noindent
We watch it climb: integers $\to$ \emph{a field has no exotic fixed points} $\to$
\emph{Gold is APN} $\to$ \emph{Kasami, even $k$, edge case}.\\[2pt]
Each floor is short. Read each as \emph{math} $\;\leftrightarrow\;$ \emph{Lean} $\;\leftrightarrow\;$
\emph{a typing judgement}.
\end{abstract}

\bigskip

% =====================================================================
\section*{The ground floor: a number}

\noindent
A \textbf{Mersenne number} is $2^a-1$. The fact:
\[
\gcd\bigl(2^a-1,\;2^b-1\bigr)\;=\;2^{\gcd(a,b)}-1 .
\]

\begin{lstlisting}
theorem gcd_mersenne (m n : ℕ) :
    Nat.gcd (2 ^ m - 1) (2 ^ n - 1) = 2 ^ (Nat.gcd m n) - 1 :=
  pow_sub_one_gcd_pow_sub_one 2 m n
\end{lstlisting}

\note{It is a \emph{term} proof: the right-hand side \emph{is} the proof.
Mathlib already carries this under the name \texttt{pow\_sub\_one\_gcd\_pow\_sub\_one}.
Our file only renames it and reads it slowly.}

\bigskip
\subsection*{Why it is true (the picture)}

\noindent
Think of the cyclic group $\mathbb{Z}/(2^a-1)$. The integer $2$ acts by
multiplication, and its order is exactly $a$ (because $2^a\equiv 1$ and no
smaller power does). So
\[
2^a\equiv 1 \pmod{N}
\quad\Longleftrightarrow\quad
(2^a-1)\mid N .
\]
Now $d\mid 2^a-1$ and $d\mid 2^b-1$ together say $2^a\equiv 2^b\equiv 1$ in the
group of units mod $d$; the smallest such exponent is $\gcd(a,b)$. Hence the
common factors are exactly $2^{\gcd(a,b)}-1$.

\[
\begin{tikzcd}[column sep=large]
2^a-1 \arrow[dr, "\gcd"'] & & 2^b-1 \arrow[dl, "\gcd"] \\
& 2^{\gcd(a,b)}-1 &
\end{tikzcd}
\]

\note{The map $a\mapsto 2^a-1$ turns $\gcd$ of \emph{exponents} into $\gcd$ of
\emph{Mersenne numbers}. It is a tiny ``functor'': it sends the lattice
$(\mathbb{N},\gcd)$ into the lattice $(\mathbb{N},\gcd)$, monotonically.}

\bigskip

% =====================================================================
\section{First floor: a field has no exotic fixed points}

\noindent
Move to $F=\mathrm{GF}(2^n)$, a field of size $2^n$. Ask: which $x$ satisfy
$x^{2^k}=x$? In char $2$ this is the Frobenius fixed-point equation.

\[
\Gamma=\bigl(F:\mathrm{Type},\;[\mathrm{Field}\;F],\;[\mathrm{Fintype}\;F],\;[\mathrm{CharP}\;F\;2]\bigr)
\]

\begin{lstlisting}
lemma frob_fixed_implies_GF2 {F : Type*} [Field F] [Fintype F] [CharP F 2]
    {n : ℕ} (hn : Fintype.card F = 2 ^ n)
    {k : ℕ} (hk : 0 < k) (hcop : Nat.Coprime k n)
    {x : F} (hfixed : x ^ (2 ^ k) = x) :
    x = 0 ∨ x = 1
\end{lstlisting}

\bigskip
\subsection*{Where the gcd enters}

\noindent
Suppose $x\neq 0$. Two facts about its multiplicative order:
\begin{itemize}\itemsep4pt
  \item from $x^{2^k}=x$ we get $x^{2^k-1}=1$;
  \item from finiteness (Fermat) we get $x^{2^n-1}=1$.
\end{itemize}
So the order of $x$ divides \emph{both} $2^k-1$ and $2^n-1$, hence divides their
gcd:
\[
x^{\gcd(2^k-1,\,2^n-1)} = 1 .
\]
Now apply the ground floor:
\[
\gcd(2^k-1,\,2^n-1)=2^{\gcd(k,n)}-1 = 2^{1}-1 = 1
\qquad(\text{since }\gcd(k,n)=1).
\]
Therefore $x^{1}=1$, i.e. $x=1$. So $x\in\{0,1\}$.

\note{The hypothesis \texttt{Nat.Coprime k n} ($\gcd(k,n)=1$) is the \emph{only}
reason the answer is the prime subfield $\{0,1\}=\mathbb{F}_2$ and nothing larger.
Drop coprimality and the fixed field grows to $\mathrm{GF}(2^{\gcd(k,n)})$.}

\[
\inferrule*[Right=\textsc{order}]
{
  \;x^{2^k-1}=1\;
  \\
  \;x^{2^n-1}=1\;
  \\
  \gcd(2^k-1,2^n-1)\overset{\text{ground}}{=}2^{\gcd(k,n)}-1=1
}
{\; x^{1}=1 \;}
\]

\bigskip

% =====================================================================
\section{Second floor: Gold is APN}

\noindent
The \textbf{Gold} function is $x\mapsto x^{2^k+1}$. \textbf{APN} (almost perfect
nonlinear) means: for every $a\neq 0$, the equation
$f(x+a)+f(x)=b$ has at most $2$ solutions.

\begin{lstlisting}
theorem gold_is_apn {F : Type*} [Field F] [Fintype F] [CharP F 2]
    {n : ℕ} (hn : Fintype.card F = 2 ^ n)
    (k : ℕ) (hk : 0 < k) (hcop : Nat.Coprime k n) :
    IsAPN (fun x : F => x ^ (2 ^ k + 1))
\end{lstlisting}

\bigskip
\subsection*{The collision collapses to a fixed point}

\noindent
Expand the Gold difference (char $2$, Freshman's Dream):
\[
(x+a)^{2^k+1}+x^{2^k+1}
= a^{2^k}x + a\,x^{2^k} + a^{2^k+1}.
\]
A collision of $x$ and $y$ forces (after cancelling) the homogeneous equation
\[
a^{2^k}(x+y) + a(x+y)^{2^k}=0 .
\]
Divide by $a^{2^k+1}$ and set $u=(x+y)/a$:
\[
u^{2^k}=u .
\]
That is the first-floor equation! By \texttt{frob\_fixed\_implies\_GF2},
$u\in\{0,1\}$, i.e. $y=x$ or $y=x+a$. At most two solutions. APN. \hfill$\blacksquare$

\note{In Lean this is literally one call:
\texttt{frob\_fixed\_implies\_GF2 hn hk hcop} closes the kernel step inside
\texttt{gold\_is\_apn}. The gcd fact is doing the load-bearing work two floors down.}

\[
\begin{tikzcd}[row sep=small, column sep=large]
\text{collision} \arrow[r, Rightarrow] & u^{2^k}=u \arrow[r, Rightarrow, "{\texttt{frob\_fixed}}"] & u\in\{0,1\} \arrow[r, Rightarrow] & \text{APN}
\end{tikzcd}
\]

\bigskip

% =====================================================================
\section{Third floor: Kasami, even $k$, at the edge}

\noindent
The Kasami exponent is $d_k = 2^{2k}-2^k+1$. The main theorem
\texttt{kasami\_is\_apn} historically needed \textbf{$k$ odd}. The even-$k$ case
is handled by a Frobenius twist that swaps $k\leftrightarrow n-k$. At the
\emph{edge} $n-k=1$ the companion exponent is the Gold exponent $d_1=3=2^1+1$:

\begin{lstlisting}
lemma kasami_one_is_apn {F : Type*} [Field F] [Fintype F] [CharP F 2]
    {n : ℕ} (hn : Fintype.card F = 2 ^ n)
    (hcop : Nat.Coprime 1 n) :
    IsAPN (fun x : F => x ^ kasamiExp 1) := by
  rw [kasamiExp_one]      -- kasamiExp 1 = 3
  exact gold_is_apn hn 1 one_pos hcop
\end{lstlisting}

\note{So the even-$k$ Kasami argument \emph{falls back on Gold} at its boundary,
and Gold falls back on the fixed-point lemma, which falls back on the Mersenne
gcd. Four floors, one staircase.}

\bigskip

% =====================================================================
\section*{The staircase}

\[
\begin{tikzcd}[row sep=large, column sep=huge]
\boxed{\gcd(2^a{-}1,2^b{-}1)=2^{\gcd(a,b)}{-}1}
  \arrow[d, Rightarrow, "{\text{order argument}}"] \\
\texttt{frob\_fixed\_implies\_GF2}\;:\; x^{2^k}=x \Rightarrow x\in\{0,1\}
  \arrow[d, Rightarrow, "{\text{collision} \to \text{kernel}}"] \\
\texttt{gold\_is\_apn}
  \arrow[d, Rightarrow, "{\text{edge } n-k=1}"] \\
\texttt{kasami\_is\_apn\_general}\;\;(\text{no parity restriction})
\end{tikzcd}
\]

\bigskip
\subsection*{One-screen summary}

\begin{center}
\renewcommand{\arraystretch}{1.6}
\begin{tabular}{l l l}
\textbf{floor} & \textbf{statement} & \textbf{engine} \\
\hline
ground & $\gcd(2^a{-}1,2^b{-}1)=2^{\gcd(a,b)}{-}1$ & cyclic order of $2$ \\
$1$ & $x^{2^k}{=}x \Rightarrow x\in\{0,1\}$ & order $\mid \gcd = 1$ \\
$2$ & Gold $x^{2^k+1}$ is APN & collision $\to$ fixed point \\
$3$ & Kasami even-$k$ edge & companion $= $ Gold ($d_1=3$) \\
\end{tabular}
\end{center}

\bigskip
\noindent
\colorbox{leanbg}{\parbox{\dimexpr\linewidth-2\fboxsep}{\centering
\texttt{gcd\_mersenne} $\;\Rightarrow\;$ \texttt{frob\_fixed\_implies\_GF2}
$\;\Rightarrow\;$ \texttt{gold\_is\_apn} $\;\Rightarrow\;$ \texttt{kasami\_is\_apn\_general}.}}

\end{document}
